Everything is computed in Gauge Theory Gravity (Definition \claimref{obj.synthesis.gtg-flat-spacetime-gauge-fields}). The Schwarzschild field of Fact \claimref{pairing.imported-schwarzschild-connection-curvature} in the Painlev\'e-Gullstrand gauge $\hb = 1 - \beta P$, $P(a) = (a\cdot e_{t})e_{r}$, $\beta = \sqrt{2M/r}$, has, on a constant orthonormal physical frame $\{e_{a}\}$, $\hb^{-1}(e_{t}) = e_{0} + \beta e_{1}$, $\hb^{-1}(e_{r}) = e_{1}$, $\hb^{-1}(e_{\vartheta}) = r\,e_{2}$, $\hb^{-1}(e_{\varphi}) = r\sin\vartheta\,e_{3}$ ($e_{t}, \dots, e_{\varphi}$ the background frame of the labels); the coframe of the statement is the set of components of $\hb^{-1}(dx)$. $\Omega(e_{c})$ means $\Omega(\hb(e_{c}))$.

\emph{The gauge field from the torsion equation.} The torsion equation of the definition, read on the constant physical frame, is $\hb^{T}(\nabla\wedge f^{a}) = \sum_{b} e^{b}\wedge(\Omega(e_{b})\cdot e^{a})$ with $\hb^{T}(f^{a}) = e^{a}$; here $f^{0} = \nabla T$, $f^{1} = \nabla r + \beta\nabla T$, $f^{2} = r\nabla\vartheta$, $f^{3} = r\sin\vartheta\,\nabla\varphi$, and $\nabla\wedge f^{1} = \beta'\,\nabla r\wedge\nabla T$ with $\beta' = d\beta/dr = -\beta/2r = -\beta^{3}/4M$. The equation is linear and injective in $\Omega$ (twenty-four equations for twenty-four components: the uniqueness of the Levi-Civita connection in frame form), and its solution is
\[
\Omega(e_{0}) = 0, \quad \Omega(e_{1}) = -\beta'\,e_{1}e_{0}, \quad \Omega(e_{2}) = \frac{(\beta e_{0} + e_{1})\wedge e_{2}}{r}, \quad \Omega(e_{3}) = \frac{(\beta e_{0} + e_{1})\wedge e_{3}}{r} + \frac{\cot\vartheta}{r}\,e_{2}\wedge e_{3}.
\]
The check solves the torsion equation exactly from the coframe alone at five rational points outside, on and inside the horizon locus ($\beta = \tfrac45, 1, 2, \tfrac13, \tfrac35$), finds rank $24$ and these bivectors. Their components $\omega^{a}{}_{b}(e_{c}) = e^{a}\cdot(\Omega(e_{c})\cdot e_{b})$ are the connection forms of the statement: $(e_{1}\cdot D)e_{1} = \Omega(e_{1})\cdot e_{1} = -\beta' e_{0}$, $(e_{2}\cdot D)e_{1} = e_{2}/r$, $(e_{2}\cdot D)e_{2} = -(\beta e_{0} + e_{1})/r$, $(e_{3}\cdot D)e_{2} = (\cot\vartheta/r)e_{3}$.

\emph{The classes.} The pairing class of the outgoing vacuum $\ell = e_{0} + e_{1}$ is $\pv(B) = \ip{B\cdot\ell}{m}$ with $m = e_{2} + I_{5}e_{3}$ in the odd extension $\Cl(2,3)$, $I_{5} = I_{4}e$ central of square $-1$ (\claimref{std-algebra.complexification-is-the-odd-extension}), the scalar product extended over $\mathrm{span}\{1, I_{5}\}$ being the symmetric product $\tfrac12(vm + mv)$. Since $(u\wedge e_{2})\cdot\ell = -(u\cdot\ell)\,e_{2}$ for $u = \beta e_{0} + e_{1}$, with $u\cdot\ell = \beta - 1$,
\[
\Omega(e_{2})\cdot\ell = \frac{1 - \beta}{r}\,e_{2}, \qquad \Omega(e_{3})\cdot\ell = \frac{1 - \beta}{r}\,e_{3},
\]
so $\pv(\Omega(e_{2})) = \tfrac{1 - \beta}{r}\ip{e_{2}}{m} = -\tfrac{1 - \beta}{r}$ and $\pv(\Omega(e_{3})) = -I_{5}\tfrac{1 - \beta}{r}$, while $\Omega(e_{0}) = 0$ and $\Omega(e_{1})\cdot\ell = -\beta'\,\ell$ give class $0$. These vanish exactly where $\beta = 1$ and have the sign of $\beta - 1$. The screen trace of $A \mapsto (A\cdot D)\ell$ is $-\sum_{c = 2,3} e_{c}\cdot(\Omega(e_{c})\cdot\ell) = 2(1 - \beta)/r$, the expansion of the outgoing congruence (by Fact \claimref{pairing.imported-spin-coefficients-null-tetrad} the class on $\mb$, $-2(1 - \beta)/r$, is the spin coefficient $\rho$); it vanishes on the horizon locus and is negative in the trapped region. Along $\ell$, $(\ell\cdot D)\ell = (\Omega(e_{0}) + \Omega(e_{1}))\cdot\ell = -\beta'\,\ell = (\beta^{3}/4M)\,\ell$. For the ingoing bit $n = e_{0} - e_{1}$, $u\cdot n = \beta + 1$, so $\Omega(e_{2})\cdot n = -\tfrac{1 + \beta}{r}e_{2}$ and the classes of $\operatorname{span}\{n, m\}$ are $+\tfrac{1 + \beta}{r}$ and $+I_{5}\tfrac{1 + \beta}{r}$, positive for every $\beta > 0$: at $M = \tfrac12$ and $\beta = \tfrac45, 1, 2$ (so $r = 2M/\beta^{2} = \tfrac{25}{16}, 1, \tfrac14$) they are $\tfrac{144}{125}$, $2$ and $12$. The check computes every class in $\Cl(2,3)$ with $I_{5}$ from the solved gauge field.

\emph{Covariance with the static frame.} The static frame is built inside the same gauge, $e_{0}^{s} = (e_{0} + \beta e_{1})/F$, $e_{1}^{s} = (e_{1} + \beta e_{0})/F$, $F = \sqrt{1 - \beta^{2}}$, and shares $e_{2}, e_{3}$; its null vector is $\ell_{s} = e_{0}^{s} + e_{1}^{s} = \tfrac{1 + \beta}{F}\,\ell$, so $\ell = \tfrac{F}{1 + \beta}\ell_{s}$, the rescaling of Proposition \claimref{pairing.two-frames-share-vacuum-selected-transport}(i). The factor depends on $r$ alone and $\hb(e_{2}), \hb(e_{3})$ have no radial component, so $(e_{c}\cdot D)\ell = \tfrac{F}{1 + \beta}(e_{c}\cdot D)\ell_{s}$ for $c = 2, 3$, and the scaling law of Proposition \claimref{pairing.pairing-class-transverse-derivative-null-direction} gives $\pv_{\mathrm{PG}} = \tfrac{F}{1 + \beta}\pv_{\mathrm{static}}$ in both transverse directions. At $(F, \beta) = (\tfrac35, \tfrac45)$, $M = \tfrac12$: $\pv_{\mathrm{PG}} = -\tfrac{1/5}{25/16} = -\tfrac{16}{125}$ and $\pv_{\mathrm{static}} = -F/r = -\tfrac{48}{125}$, with factor $\tfrac13$.

\emph{The curvature class.} The field strength of the solved gauge field, $\mathcal{R}(B) = -R(\hb(B))$ on the physical frame planes, is $2M/r^{3}$ times the plane on $e_{0}\wedge e_{1}$ and $e_{2}\wedge e_{3}$ and $-M/r^{3}$ on the four mixed planes, the same coefficients as in the static frame (the check computes it from the solved field at the five points). The class of the loop $(n, e_{2})$ is $\pv(-\mathcal{R}(n\wedge e_{2})) = 2M/r^{3}$, equal to $1$ on the horizon locus at $M = \tfrac12$, and it is unchanged by the radial boost between the frames because $\ell$ and $n$ scale inversely. So the curvature classes do not vanish where the connection classes do.

\emph{In components} (a labelled reading). The torsion equation is the first structure equation $d\theta^{a} + \omega^{a}{}_{b}\wedge\theta^{b} = 0$ of the coframe $\theta^{a}$: $d\theta^{1} = \beta'\,dr\wedge dT = \beta'\,\theta^{1}\wedge\theta^{0}$ fixes $\omega^{0}{}_{1} = -\beta'\theta^{1}$, and $d\theta^{2} = (\theta^{1} - \beta\theta^{0})\wedge\theta^{2}/r$ gives $\omega^{2}{}_{1}$ and $\omega^{2}{}_{0}$, the $\hat3$ forms following by spherical symmetry; $\nabla_{e_{2}}\ell = (\omega^{2}{}_{0} + \omega^{2}{}_{1})(e_{2})\,e_{2}$ is the first display above. The earlier component checks are kept beside the GTG ones.
